\clearpage
\begin{explanationbox}
\textbf{先问三件事：为什么要分类？为什么分界是 $b/a$？为什么有两个不同最小值？}

\textbf{第一件：先画渐近线，再转动过焦点的直线。}
双曲线的两条渐近线斜率是 $\pm b/a$；它们不仅决定开口的方向，
也决定过焦点的直线究竟先穿过\textbf{左右两支}，还是\textbf{右支两次}。
\begin{itemize}
\item $|k|<b/a$（较平）：两个交点在左右不同支，称\textbf{异支弦}；
\item $|k|>b/a$（较陡）：两个交点同在右支，称\textbf{同支弦}；
\item $k=\pm b/a$（临界）：直线平行渐近线，只剩一个有限交点；
\item 竖直的 $x=c$ 不是斜率有限的 $k$，它是合法的\textbf{同支弦}。
\end{itemize}
对于左焦点 $F_1$，对应为\textbf{左支}两次，其他结论不变。
\begin{center}% 精确解析交点，a=3,b=4,c=5，曲线参数由标准方程求出。
% 左：k=1/4 (<4/3)，跨两支。右：k=4 (>4/3)，同右支。
% 对 k=1/4，t 是 (247/16)t^2+160t+256=0 的两实根：(-1280±192sqrt(17))/247。
% 对 k=4，t 是 -128t^2+160t+256=0 的两实根：(5±3sqrt(17))/8。
\begin{tikzpicture}[x=0.39cm,y=0.39cm,>=stealth,font=\small]
\begin{scope}
  \draw[->,gray] (-7.2,0)--(7.3,0) node[right]{$x$};
  \draw[->,gray] (0,-7.8)--(0,7.8) node[above]{$y$};
  \draw[thick,blue!75!black,domain=-7:7,samples=90] plot ({3*sqrt(1+\x*\x/16)},{\x});
  \draw[thick,blue!75!black,domain=-7:7,samples=90] plot ({-3*sqrt(1+\x*\x/16)},{\x});
  \draw[densely dashed,gray!75,domain=-7:7] plot ({\x},{4/3*\x});
  \draw[densely dashed,gray!75,domain=-7:7] plot ({\x},{-4/3*\x});
  \pgfmathsetmacro{\ta}{(-1280+192*sqrt(17))/247}
  \pgfmathsetmacro{\tb}{(-1280-192*sqrt(17))/247}
  \draw[dashed,red!75!black] (5,0)--({5+\tb},{\tb/4});
  \draw[very thick,red!75!black] ({5+\ta},{\ta/4})--({5+\tb},{\tb/4});
  \fill[red!75!black] ({5+\ta},{\ta/4}) circle (1.8pt) node[below left=5pt]{$B$};
  \fill[red!75!black] ({5+\tb},{\tb/4}) circle (1.8pt) node[left=7pt]{$A$};
  \fill (5,0) circle (2pt) node[above right]{$F_2$};
  \node at (0,-9.0) {$k=1/4$: 异支};
\end{scope}
\begin{scope}[xshift=6.6cm]
  \draw[->,gray] (-7.2,0)--(8.0,0) node[right]{$x$};
  \draw[->,gray] (0,-9.6)--(0,9.8) node[above]{$y$};
  \draw[thick,blue!75!black,domain=-9:9,samples=110] plot ({3*sqrt(1+\x*\x/16)},{\x});
  \draw[thick,blue!75!black,domain=-9:9,samples=110] plot ({-3*sqrt(1+\x*\x/16)},{\x});
  \draw[densely dashed,gray!75,domain=-7:7] plot ({\x},{4/3*\x});
  \draw[densely dashed,gray!75,domain=-7:7] plot ({\x},{-4/3*\x});
  \pgfmathsetmacro{\tc}{(5+3*sqrt(17))/8}
  \pgfmathsetmacro{\td}{(5-3*sqrt(17))/8}
  \draw[very thick,red!75!black] ({5+\tc},{4*\tc})--({5+\td},{4*\td});
  \fill[red!75!black] ({5+\tc},{4*\tc}) circle (1.8pt) node[above left=6pt]{$A$};
  \fill[red!75!black] ({5+\td},{4*\td}) circle (1.8pt) node[below left=5pt]{$B$};
  \fill (5,0) circle (2pt) node[below right]{$F_2$};
  \node at (0,-11) {$k=4$: 同右支};
\end{scope}
\end{tikzpicture}
\end{center}

\tcbbreak
\textbf{第二件：把几何分类放到弦长变化图上。}
\begin{center}% Actual L(theta), a=3,b=4,c=5; horizontal degree/10, vertical .14L.
\begin{tikzpicture}[font=\small,>=stealth]
\draw[gray!45,rounded corners=2pt] (0,0) rectangle (14.2,9.2);
\node at (7.1,8.55) {转动直线看弦长：临界处无界增长，水平和竖直分别取最小值};
\node at (7.1,7.75) {图取 $a=3,b=4,c=5$，只看 $0\le\theta\le\pi/2$};
\begin{scope}[shift={(2.3,1.6)}]
\draw[gray,->] (0,0)--(9.7,0) node[right] {$\theta$};
\draw[gray,->] (0,0)--(0,5.8) node[above] {$L$};
\draw[gray,dashed] ({acos(3/5)/10},0)--({acos(3/5)/10},5.8);
\draw[red!75!black,very thick,->,domain=0:{acos(sqrt((9+96/40)/25))},samples=100,variable=\v]
 plot ({\v/10},{.14*96/(25*cos(\v)^2-9)});
\draw[purple!80!black,very thick,<-,domain={acos(sqrt((9-96/40)/25))}:90,samples=100,variable=\v]
 plot ({\v/10},{.14*96/(9-25*cos(\v)^2)});
\fill[red!75!black] (0,.84) circle(2.2pt) node[left=6pt] {$2a=6$};
\fill[purple!80!black] (9,{.14*32/3}) circle(2.2pt) node[above right=5pt] {$2b^2/a=32/3$};
\node[below=5pt] at (0,0) {$0$};
\node[below=5pt] at (9,0) {$\pi/2$};
\node[below=5pt] at ({acos(3/5)/10},0) {$\theta_0=\arccos(3/5)$};
\node[red!75!black] at (2,2.8) {异支};
\node[purple!80!black] at (7.4,3.4) {同支};
\end{scope}
\node at (7.1,.45) {临界角 $\theta_0$ 对应渐近线方向；实心点说明两个最小值都能取到。};
\end{tikzpicture}
\end{center}
两侧曲线都向临界方向升高；水平弦和竖直弦是两类范围各自的最低点。
\textbf{不要被绝对值吓到，它只表示分母变号。}
设 $u=k^2$，临界值 $u_0=b^2/a^2$。只须记住
\[
\begin{array}{ll}
0\le u<u_0:& \displaystyle L=\frac{2ab^2(1+u)}{b^2-a^2u}\quad(\text{异支}),\\[5pt]
u>u_0:& \displaystyle L=\frac{2ab^2(1+u)}{a^2u-b^2}\quad(\text{同支}).
\end{array}
\]
\textbf{第三件：把「最小值」变成「比最小值多多少」。}
异支时
\[
\boxed{L-2a=\frac{2ac^2u}{b^2-a^2u}\ge0}
\]
只在 $u=0$ 取等，所以水平弦最短。
同支、斜率有限时
\[
\boxed{L-\frac{2b^2}{a}=
\frac{2b^2c^2}{a(a^2u-b^2)}>0.}
\]
有限斜率不可能取等；必须单独检查竖直线 $x=c$：
与曲线交于 $\bigl(c,\pm b^2/a\bigr)$，其长\textbf{恰为} $2b^2/a$。

\textbf{极限直觉：}转向平行渐近线时，远处那个交点跑向无限远，
所以弦长发散；转到竖直时，两个交点关于 $x$ 轴对称，得到同支最短弦。
\end{explanationbox}
